\notechapter{N-20221221}{Completion, Uniform Structure, and Quotient Topology}

\section{A Cauchy sequence can describe a point missing from the space}

A sequence \((x_n)\) in a metric space \((X,d)\) is Cauchy when for every \(\varepsilon>0\) there is \(N\) such that
\[
  d(x_m,x_n)<\varepsilon
  \qquad(m,n\ge N).
\]
It records points that eventually become mutually indistinguishable at every metric scale.  A complete space contains a limit for every such sequence.

The rational sequence
\[
  q_n=\frac{\lfloor10^n\sqrt2\rfloor}{10^n}
\]
is Cauchy in \(\mathbb Q\), but it has no rational limit.  Nothing is wrong with the sequence; the intended limit is absent from the space.  Completion constructs a space in which every coherent approximation process has a point to name.

Two Cauchy sequences should represent the same new point when their mutual distance tends to zero:
\[
  (x_n)\sim(y_n)
  \quad\Longleftrightarrow\quad
  d(x_n,y_n)\longrightarrow0.
\]
Reflexivity and symmetry are immediate.  If \((x_n)\sim(y_n)\) and \((y_n)\sim(z_n)\), then
\[
  d(x_n,z_n)
  \le d(x_n,y_n)+d(y_n,z_n)\longrightarrow0,
\]
so the relation is transitive.

\section{The completion metric is a limit of mutual distances}

Let \(\widehat X\) be the set of equivalence classes of Cauchy sequences.  Define
\[
  \widehat d([(x_n)],[(y_n)])
  =\lim_{n\to\infty}d(x_n,y_n).
\]
The scalar sequence on the right is Cauchy because
\[
\begin{aligned}
  |d(x_n,y_n)-d(x_m,y_m)|
  &\le d(x_n,x_m)+d(y_n,y_m).
\end{aligned}
\]
Hence the limit exists in \(\mathbb R\).

The value does not depend on the chosen representatives.  If
\[
  (x_n)\sim(x_n'),
  \qquad
  (y_n)\sim(y_n'),
\]
then
\[
  |d(x_n,y_n)-d(x_n',y_n')|
  \le d(x_n,x_n')+d(y_n,y_n')\longrightarrow0.
\]
Thus both representative pairs give the same limit.

Nonnegativity and symmetry pass directly to the limit.  If \(\widehat d=0\), then \(d(x_n,y_n)\to0\), so the classes are equal.  Finally,
\[
  d(x_n,z_n)
  \le d(x_n,y_n)+d(y_n,z_n)
\]
passes to the limit and gives the triangle inequality.  Therefore \(\widehat d\) is a genuine metric.

\section{The original space sits isometrically and densely inside its completion}

Send \(x\in X\) to the constant sequence:
\[
  \iota(x)=[(x,x,x,\ldots)].
\]
Then
\[
  \widehat d(\iota(x),\iota(y))=d(x,y),
\]
so \(\iota\) is an isometric embedding.

Its image is dense.  Let
\[
  \xi=[(x_n)]\in\widehat X.
\]
For fixed \(N\),
\[
  \widehat d(\xi,\iota(x_N))
  =\lim_{m\to\infty}d(x_m,x_N).
\]
Because \((x_n)\) is Cauchy, this quantity tends to zero as \(N\to\infty\).  Thus
\[
  \iota(x_N)\longrightarrow\xi.
\]
Every completed point is a limit of original points.

For \(X=\mathbb Q\), the decimal truncations \((q_n)\) of \(\sqrt2\) define a class in \(\widehat{\mathbb Q}\).  It cannot equal any constant rational class: if it equalled \(r\in\mathbb Q\), then
\[
  |q_n-r|\longrightarrow0,
\]
forcing \(r=\sqrt2\).  The new class is exactly the missing irrational point described by rational approximations.

\section{A diagonal approximation proves that the completed space is complete}

Let \((\xi_r)\) be a Cauchy sequence in \(\widehat X\).  Passing to a subsequence if necessary, arrange that
\[
  \widehat d(\xi_r,\xi_{r+1})<2^{-r}.
\]
Density of \(\iota(X)\) lets us choose \(y_r\in X\) with
\[
  \widehat d(\iota(y_r),\xi_r)<2^{-r}.
\]
Then
\[
\begin{aligned}
  d(y_r,y_{r+1})
  &\le
  \widehat d(\iota(y_r),\xi_r)
  +\widehat d(\xi_r,\xi_{r+1})\\
  &\qquad
  +\widehat d(\xi_{r+1},\iota(y_{r+1}))\\
  &<2^{-r}+2^{-r}+2^{-(r+1)}.
\end{aligned}
\]
The geometric tail shows that \((y_r)\) is Cauchy in \(X\).  Let
\[
  \xi=[(y_r)]\in\widehat X.
\]
The construction gives
\[
  \xi_r\longrightarrow\xi
\]
along the chosen subsequence, and since the original sequence was Cauchy, the whole sequence has the same limit.  Hence \(\widehat X\) is complete.

The argument exposes the mechanism: approximate each abstract completed point by an original point with rapidly shrinking error, then use one diagonal Cauchy sequence in \(X\) to encode the limit.

\section{Uniformly continuous maps extend because they preserve Cauchy information}

Let \(Y\) be complete and let
\[
  f:X\to Y
\]
be uniformly continuous.  If \((x_n)\) is Cauchy, then \((f(x_n))\) is Cauchy.  Indeed, for every \(\varepsilon>0\), choose one \(\delta>0\) working simultaneously for all points:
\[
  d_X(x,x')<\delta
  \quad\Longrightarrow\quad
  d_Y(f(x),f(x'))<\varepsilon.
\]
The Cauchy condition eventually places every pair \(x_m,x_n\) within \(\delta\).

Define
\[
  \widehat f:\widehat X\to Y,
  \qquad
  \widehat f([(x_n)])
  =\lim_{n\to\infty}f(x_n).
\]
If \((x_n)\sim(y_n)\), uniform continuity gives
\[
  d_Y(f(x_n),f(y_n))\longrightarrow0,
\]
so the definition is independent of representatives.  It extends \(f\) because constant sequences retain their values.

The extension is unique.  If \(g:\widehat X\to Y\) is continuous and agrees with \(f\) on \(X\), then for
\[
  \xi=\lim_n\iota(x_n)
\]
one has
\[
  g(\xi)
  =\lim_ng(\iota(x_n))
  =\lim_nf(x_n)
  =\widehat f(\xi).
\]
Thus completion is characterized by the extension property
\[
  \boxed{
  \text{uniformly continuous maps into complete spaces extend uniquely}.}
\]

Uniform continuity is essential.  The continuous map
\[
  f:(0,1)\to\mathbb R,
  \qquad
  f(x)=\frac1x
\]
has no continuous extension to the completion \([0,1]\), because it diverges along \(x_n=1/n\to0\).

\section{The universal property makes completion unique}

Suppose
\[
  \iota_X:X\hookrightarrow Y,
  \qquad
  j_X:X\hookrightarrow Z
\]
are dense isometric embeddings into complete metric spaces.  Apply the extension property to the uniformly continuous map
\[
  j_X\circ\iota_X^{-1}:\iota_X(X)\to Z.
\]
It extends uniquely to an isometry
\[
  F:Y\to Z.
\]
Reversing the roles gives
\[
  G:Z\to Y.
\]
Both composites agree with the identity on the dense copy of \(X\).  By uniqueness of extension,
\[
  GF=\operatorname{id}_Y,
  \qquad
  FG=\operatorname{id}_Z.
\]
Therefore any two completions are uniquely isometric in a way that fixes \(X\).

The quotient of Cauchy sequences is one concrete construction; the universal property is what makes the result canonical.

\section{Homeomorphism forgets the uniform scale of closeness}

The interval \((0,1)\) and the real line are homeomorphic via
\[
  h(x)=\tan\bigl(\pi(x-1/2)\bigr).
\]
Yet \((0,1)\) is bounded and incomplete in its Euclidean metric, while \(\mathbb R\) is unbounded and complete.  Boundedness and completeness are therefore not topological invariants.

The map is not uniformly continuous.  Take
\[
  x_n=1-\frac1n,
  \qquad
  y_n=1-\frac2n.
\]
Then
\[
  |x_n-y_n|=\frac1n\longrightarrow0,
\]
but
\[
  h(x_n)=\cot\frac\pi n,
  \qquad
  h(y_n)=\cot\frac{2\pi}n,
\]
and their difference grows on the order of \(n\).  No single \(\delta\) controls the map near both endpoints.

Topology records which sets are neighborhoods.  Uniform structure additionally records which pairs are uniformly close across the whole space.  Completion depends on the latter, because Cauchy sequences are uniform objects.

\section{A continuous bijection needs extra hypotheses to be a homeomorphism}

The map
\[
  f:[0,2\pi)\to S^1,
  \qquad
  f(t)=(\cos t,\sin t)
\]
is continuous and bijective, but its inverse is not continuous at \((1,0)\).  Points approaching \((1,0)\) from the lower semicircle have parameters tending to \(2\pi\), while the point itself has parameter \(0\).

A useful positive theorem is the compact--Hausdorff criterion.

\begin{theorem}
If \(X\) is compact, \(Y\) is Hausdorff, and \(f:X\to Y\) is a continuous bijection, then \(f\) is a homeomorphism.
\end{theorem}

\begin{proof}
A closed subset \(C\subseteq X\) is compact.  Its image \(f(C)\) is compact in \(Y\), hence closed because \(Y\) is Hausdorff.  Thus \(f\) is a closed map.  For any closed \(C\subseteq X\),
\[
  (f^{-1})^{-1}(C)=f(C)
\]
is closed, so \(f^{-1}\) is continuous.
\end{proof}

The failed half-open-interval example lacks compactness.  The theorem identifies exactly which global hypotheses make inverse continuity automatic.

\section{Quotient topology is designed to control maps out of a gluing}

Let \(q:X\to Y\) be a surjection.  The quotient topology on \(Y\) is defined by
\[
  U\subseteq Y\text{ open}
  \quad\Longleftrightarrow\quad
  q^{-1}(U)\text{ open in }X.
\]
It is the strongest topology on \(Y\) making \(q\) continuous.

Its universal property is the factorization test:
\[
  g:Y\to Z\text{ is continuous}
  \quad\Longleftrightarrow\quad
  g\circ q:X\to Z\text{ is continuous}.
\]
Thus a map out of a quotient is constructed upstairs, where the original points are still visible, provided the map is constant on equivalence classes.

A subset \(A\subseteq X\) is saturated when it is a union of equivalence classes:
\[
  A=q^{-1}(q(A)).
\]
For arbitrary \(A\),
\[
  q(A)\text{ is open in }Y
  \quad\Longleftrightarrow\quad
  q^{-1}(q(A))\text{ is open in }X.
\]
A quotient map need not send every open set to an open set.  On \([0,1]\), collapse the entire interval \([0,1/2]\) to one point.  The open set \((0,1/2)\) has saturation \([0,1/2]\), which is not open, so its image is not open in the quotient.

The quotient topology controls precisely the saturated information that survives the identification.

\section{Endpoint gluing becomes the circle by quotient factorization}

Let
\[
  q:[0,1]\to[0,1]/(0\sim1)
\]
be the endpoint quotient, and define
\[
  F(t)=e^{2\pi it}.
\]
Since
\[
  F(0)=F(1),
\]
the map is constant on equivalence classes.  The quotient universal property produces a unique continuous map
\[
  \overline F:[0,1]/(0\sim1)\to S^1
\]
with
\[
  \overline F\circ q=F.
\]
It is bijective: every point of the circle has a parameter in \([0,1]\), and the only repeated parameter values are the identified endpoints.

The quotient of the compact interval is compact, while \(S^1\) is Hausdorff.  By the compact--Hausdorff theorem,
\[
  \boxed{
  [0,1]/(0\sim1)\cong S^1.}
\]

A neighborhood of the glued point has preimage
\[
  [0,\varepsilon)\cup(1-\varepsilon,1],
\]
so two one-sided endpoint neighborhoods become one two-sided circular neighborhood.  The concrete local picture, the quotient factorization property, and the compactness theorem all describe the same gluing.